 \documentclass[12pt, a4paper]{article}

\usepackage[utf8]{inputenc}
\usepackage[T1]{fontenc}
\usepackage[italian]{babel}
\usepackage{geometry}
\usepackage{graphicx}
\usepackage{amsmath}
\usepackage{amssymb}
\usepackage{listings}
\usepackage{xcolor}
\usepackage{hyperref}
\usepackage{booktabs}
\usepackage{float}
\usepackage{fancyhdr}
\usepackage{titlesec}
\usepackage{enumitem}
\usepackage{microtype}

\geometry{
    top=2.5cm,
    bottom=2.5cm,
    left=3cm,
    right=3cm
}

% Stile intestazione
\pagestyle{fancy}
\fancyhf{}
\rhead{TLN - Lab 4}
\lhead{Università degli Studi di Torino}
\rfoot{\thepage}

% Colori per il codice
\definecolor{codegreen}{rgb}{0,0.6,0}
\definecolor{codegray}{rgb}{0.5,0.5,0.5}
\definecolor{codepurple}{rgb}{0.58,0,0.82}
\definecolor{backcolour}{rgb}{0.97,0.97,0.97}
\definecolor{codered}{rgb}{0.8,0,0}

% Stile codice Python
\lstdefinestyle{pythonstyle}{
    backgroundcolor=\color{backcolour},
    commentstyle=\color{codegreen},
    keywordstyle=\color{codepurple}\bfseries,
    numberstyle=\tiny\color{codegray},
    stringstyle=\color{codered},
    basicstyle=\ttfamily\footnotesize,
    breakatwhitespace=false,
    breaklines=true,
    captionpos=b,
    keepspaces=true,
    numbers=left,
    numbersep=5pt,
    showspaces=false,
    showstringspaces=false,
    showtabs=false,
    tabsize=2,
    frame=single,
    rulecolor=\color{codegray}
}

\lstset{style=pythonstyle}

% Hyperref setup
\hypersetup{
    colorlinks=true,
    linkcolor=blue,
    filecolor=magenta,
    urlcolor=cyan,
}

% ============================================================
\begin{document}
% ============================================================

% Pagina titolo
\begin{titlepage}
    \centering
    \vspace*{1cm}

    {\Large \textbf{Università degli Studi di Torino}}\\
    \vspace{0.3cm}
    {\large Dipartimento di Informatica}\\
    \vspace{0.3cm}
    {\large Corso di Laurea Magistrale in Informatica}

    \vspace{2cm}

    \rule{\linewidth}{0.5mm}\\
    \vspace{0.5cm}
    {\LARGE \textbf{Tecnologie del Linguaggio Naturale}}\\
    \vspace{0.3cm}
    {\Large \textbf{Lab 4 --- Modellazione del Linguaggio con N-gram}}\\
    \vspace{0.3cm}
    {\large \textit{Modeling social media and literary language with N-grams}}\\
    \vspace{0.5cm}
    \rule{\linewidth}{0.5mm}

    \vspace{2cm}

    \begin{tabular}{ll}
        \textbf{Studente:}  & Alessandro Olivero \\
        \textbf{Matricola:} & [INSERIRE MATRICOLA] \\
        \textbf{Docente:}   & Prof. Daniele Radicioni \\
        \textbf{A.A.:}      & 2024/2025 \\
    \end{tabular}

    \vfill
    {\large Torino, \today}
\end{titlepage}

% Indice
\tableofcontents
\newpage

% ============================================================
\section{Introduzione}
% ============================================================

Questo laboratorio ha come obiettivo la modellazione statistica del linguaggio
attraverso modelli N-gram applicati a due domini linguistici radicalmente diversi:
i \textbf{social media} (Twitter) e il \textbf{linguaggio letterario} (romanzi
dell'Ottocento).

L'ipotesi di partenza è che i due domini siano statisticamente distinguibili:
il linguaggio dei social media è caratterizzato da frasi brevi, abbreviazioni,
slang ed espressioni emotive, mentre il linguaggio letterario si distingue per
strutture sintattiche complesse, vocabolario formale e coerenza narrativa.

\subsection{Obiettivi}

\begin{itemize}
    \item Costruire modelli N-gram (unigrammi, bigrammi, trigrammi) su entrambi
          i corpora e analizzare le sequenze più frequenti.
    \item Addestrare modelli di linguaggio MLE e Laplace e confrontarne i limiti.
    \item Dimostrare le differenze tra i due domini attraverso la perplexity
          e la log-likelihood.
    \item Implementare un classificatore basato su log-likelihood capace di
          distinguere automaticamente un testo Twitter da uno letterario.
    \item Analizzare le differenze linguistiche attraverso metriche quantitative
          come il Type-Token Ratio (TTR) e la lunghezza media delle frasi.
\end{itemize}

\subsection{Strumenti utilizzati}

\begin{itemize}
    \item \textbf{Python 3.x}
    \item \textbf{NLTK} (Natural Language Toolkit): tokenizzazione, corpora,
          modelli di linguaggio
    \item \textbf{Matplotlib}: visualizzazione dei risultati
    \item \textbf{scikit-learn}: valutazione formale del classificatore
\end{itemize}

% ============================================================
\section{Dataset}
% ============================================================

\subsection{Twitter Samples (NLTK)}

Il corpus Twitter è stato ottenuto dal pacchetto \texttt{nltk.corpus.twitter\_samples},
che contiene tweet pre-classificati in positivi e negativi raccolti nel 2015.

\begin{itemize}
    \item \texttt{positive\_tweets.json}: 5.000 tweet con emoticon positive \texttt{:)}
    \item \texttt{negative\_tweets.json}: 5.000 tweet con emoticon negative \texttt{:(}
    \item \textbf{Totale}: 10.000 tweet
\end{itemize}

\begin{lstlisting}[language=Python, caption=Caricamento dataset Twitter]
from nltk.corpus import twitter_samples

tweets_pos = twitter_samples.strings('positive_tweets.json')
tweets_neg = twitter_samples.strings('negative_tweets.json')
tweets_raw = tweets_pos + tweets_neg  # 10000 tweet totali
\end{lstlisting}

Esempi di tweet dal dataset:

\begin{itemize}
    \item \textit{Positivo}: ``\#FollowFriday @France\_Inte for being top engaged members :)''
    \item \textit{Negativo}: ``hopeless for tmr :( ''
\end{itemize}

\subsection{Corpus Letterario (Gutenberg)}

Il testo letterario è stato estratto dal corpus Gutenberg di NLTK, che raccoglie
opere classiche di dominio pubblico. È stato scelto il romanzo
\textbf{\textit{Emma} di Jane Austen (1816)} come rappresentante del linguaggio
letterario formale ottocentesco.

\begin{lstlisting}[language=Python, caption=Caricamento corpus letterario]
from nltk.corpus import gutenberg

testo_letterario_raw = gutenberg.raw('austen-emma.txt')
# 887.071 caratteri totali
\end{lstlisting}

\subsection{Statistiche comparative}

\begin{table}[H]
\centering
\begin{tabular}{lrr}
\toprule
\textbf{Metrica} & \textbf{Twitter} & \textbf{Letterario} \\
\midrule
Testi totali         & 10.000 tweet    & 1 romanzo \\
Token totali         & 94.706          & 158.267 \\
Parole uniche        & 12.024          & 9.305 \\
Lunghezza media      & 16,0 parole     & 25,6 parole \\
Type-Token Ratio     & 0,1270          & 0,0588 \\
\bottomrule
\end{tabular}
\caption{Statistiche comparative dei due corpora}
\end{table}

% ============================================================
\section{Preprocessing}
% ============================================================

\subsection{Pulizia dei testi}

I due corpora richiedono procedure di pulizia diverse data la loro natura.

\subsubsection{Pulizia Twitter}

\begin{lstlisting}[language=Python, caption=Funzione di pulizia tweet]
def pulisci_tweet(testo):
    testo = re.sub(r'http\S+', '', testo)       # rimuovi URL
    testo = re.sub(r'@\w+', '', testo)           # rimuovi @menzioni
    testo = re.sub(r'#', '', testo)              # rimuovi simbolo #
    testo = re.sub(r'[^a-zA-Z\s]', '', testo)   # solo lettere
    testo = testo.lower().strip()
    return testo
\end{lstlisting}

Esempio di trasformazione:
\begin{itemize}
    \item \textbf{Prima}: \texttt{\#FollowFriday @France\_Inte for being top :)}
    \item \textbf{Dopo}: \texttt{followfriday for being top}
\end{itemize}

\subsubsection{Pulizia testo letterario}

\begin{lstlisting}[language=Python, caption=Funzione di pulizia testo letterario]
def pulisci_letterario(testo):
    testo = re.sub(r'[^a-zA-Z\s]', '', testo)
    testo = testo.lower().strip()
    return testo
\end{lstlisting}

\subsection{Tokenizzazione}

La tokenizzazione è stata effettuata con \texttt{word\_tokenize} di NLTK,
che a differenza del semplice \texttt{split()} gestisce correttamente
la punteggiatura e le contrazioni inglesi.

\begin{lstlisting}[language=Python, caption=Tokenizzazione]
from nltk.tokenize import word_tokenize

tokens_tweet = word_tokenize(' '.join(tweets_puliti))
tokens_lett  = word_tokenize(testo_lett_pulito)
\end{lstlisting}

Esempio di differenza tra \texttt{split()} e \texttt{word\_tokenize}:

\begin{table}[H]
\centering
\begin{tabular}{ll}
\toprule
\textbf{Metodo} & \textbf{Output} \\
\midrule
\texttt{split()} & \texttt{["I", "can't", "stop,"]} \\
\texttt{word\_tokenize()} & \texttt{["I", "ca", "n't", "stop", ","]} \\
\bottomrule
\end{tabular}
\caption{Confronto tra split() e word\_tokenize()}
\end{table}

\subsection{Rimozione delle Stopwords}

Le stopwords sono parole molto comuni (``the'', ``a'', ``in'') che appaiono
frequentemente in tutti i tipi di testo e quindi non sono discriminative.
Rimuoverle permette di far emergere i bigrammi veramente caratteristici
di ciascun dominio.

\begin{lstlisting}[language=Python, caption=Rimozione stopwords]
from nltk.corpus import stopwords

stop = set(stopwords.words('english'))
tokens_tweet_filtrati = [t for t in tokens_tweet if t not in stop]
tokens_lett_filtrati  = [t for t in tokens_lett  if t not in stop]
\end{lstlisting}

% ============================================================
\section{Modelli N-gram}
% ============================================================

\subsection{Cos'è un N-gram}

Un \textbf{N-gram} è una sequenza contigua di $N$ parole in un testo.
Formalmente, dato un testo $w_1, w_2, \ldots, w_T$, un N-gram è una
sequenza $(w_i, w_{i+1}, \ldots, w_{i+N-1})$.

\begin{itemize}
    \item \textbf{Unigrammi} ($N=1$): singole parole.
          Frase: ``i love cats'' $\rightarrow$ \texttt{[i, love, cats]}
    \item \textbf{Bigrammi} ($N=2$): coppie di parole consecutive.
          Frase: ``i love cats'' $\rightarrow$ \texttt{[(i,love), (love,cats)]}
    \item \textbf{Trigrammi} ($N=3$): triplette di parole consecutive.
          Frase: ``i love cats'' $\rightarrow$ \texttt{[(i,love,cats)]}
\end{itemize}

\subsection{Analisi delle frequenze}

\subsubsection{Bigrammi senza stopwords}

I bigrammi più frequenti dopo la rimozione delle stopwords sono i più
informativi per caratterizzare ciascun dominio:

\begin{table}[H]
\centering
\begin{tabular}{llr}
\toprule
\textbf{Dominio} & \textbf{Bigramma} & \textbf{Frequenza} \\
\midrule
Twitter & (wan, na)         & 123 \\
Twitter & (follow, please)  & 81 \\
Twitter & (gon, na)         & 72 \\
Twitter & (happy, birthday) & 54 \\
\midrule
Letterario & (mr, knightley)   & 248 \\
Letterario & (mrs, weston)     & 222 \\
Letterario & (mr, elton)       & 184 \\
Letterario & (miss, woodhouse) & 148 \\
\bottomrule
\end{tabular}
\caption{Top bigrammi per dominio (senza stopwords)}
\end{table}

I bigrammi Twitter riflettono il linguaggio colloquiale e le interazioni
sociali tipiche della piattaforma: \textit{wanna}, \textit{gonna},
richieste di follow reciproco, auguri di compleanno.

I bigrammi letterari sono dominati da titoli e nomi propri dei personaggi
del romanzo (Mr. Knightley, Mrs. Weston, Miss Woodhouse), riflettendo
il registro formale dell'inglese ottocentesco.

\subsubsection{Bigrammi con stopwords}

\begin{table}[H]
\centering
\begin{tabular}{llr}
\toprule
\textbf{Dominio} & \textbf{Bigramma} & \textbf{Frequenza} \\
\midrule
Twitter & (thank, you)  & 260 \\
Twitter & (i, love)     & 173 \\
Twitter & (i, miss)     & 160 \\
Twitter & (love, you)   & 155 \\
Twitter & (i, cant)     & 155 \\
\midrule
Letterario & (to, be)    & 602 \\
Letterario & (of, the)   & 563 \\
Letterario & (i, am)     & 363 \\
Letterario & (she, had)  & 322 \\
Letterario & (had, been) & 305 \\
\bottomrule
\end{tabular}
\caption{Top bigrammi per dominio (con stopwords)}
\end{table}

Osservazioni:
\begin{itemize}
    \item Twitter usa la forma contratta \textit{cant} invece di \textit{cannot},
          tipica del linguaggio informale.
    \item Il letterario usa \textit{i am} (forma completa) invece di \textit{i'm},
          tipico dell'inglese formale ottocentesco.
    \item Il letterario mostra strutture come \textit{had been} (past perfect)
          assenti nei tweet.
\end{itemize}

\subsubsection{Trigrammi}

\begin{table}[H]
\centering
\begin{tabular}{llr}
\toprule
\textbf{Dominio} & \textbf{Trigramma} & \textbf{Frequenza} \\
\midrule
Twitter & (i, love, you)        & 84 \\
Twitter & (you, so, much)       & 74 \\
Twitter & (thanks, for, the)    & 67 \\
\midrule
Letterario & (i, do, not)         & 123 \\
Letterario & (i, am, sure)        & 97 \\
Letterario & (she, could, not)    & 68 \\
Letterario & (a, great, deal)     & 63 \\
\bottomrule
\end{tabular}
\caption{Top trigrammi per dominio}
\end{table}

% ============================================================
\section{Modelli di Linguaggio}
% ============================================================

\subsection{Maximum Likelihood Estimation (MLE)}

Il modello MLE stima la probabilità di una parola dato il contesto
precedente contando le occorrenze nel corpus di training.

Per i bigrammi, la probabilità condizionale è:

\begin{equation}
P(w_i \mid w_{i-1}) = \frac{C(w_{i-1}, w_i)}{C(w_{i-1})}
\end{equation}

dove $C(w_{i-1}, w_i)$ è il numero di occorrenze del bigramma
$(w_{i-1}, w_i)$ e $C(w_{i-1})$ è il numero di occorrenze
di $w_{i-1}$.

\begin{lstlisting}[language=Python, caption=Addestramento modello MLE]
from nltk.lm import MLE
from nltk.lm.preprocessing import padded_everygram_pipeline

n = 2  # bigrammi
train_data, vocab = padded_everygram_pipeline(n, [tokens_tweet])

modello = MLE(n)
modello.fit(train_data, vocab)
\end{lstlisting}

\subsubsection{Probabilità condizionali osservate}

\begin{table}[H]
\centering
\begin{tabular}{llrr}
\toprule
\textbf{Parola} & \textbf{Contesto} &
\textbf{P (Twitter)} & \textbf{P (Letterario)} \\
\midrule
love  & i    & 0,0530 & 0,0017 \\
wait  & cant & 0,1146 & 0,0000 \\
happy & so   & 0,0138 & --- \\
mr    & said & 0,0000 & 0,0802 \\
been  & had  & ---    & 0,1885 \\
\bottomrule
\end{tabular}
\caption{Probabilità condizionali nei due modelli MLE}
\end{table}

I risultati confermano le differenze attese:
\begin{itemize}
    \item \textit{love} segue \textit{i} con probabilità 5,3\% su Twitter,
          quasi mai nel romanzo.
    \item \textit{mr} segue \textit{said} con probabilità 8\% nel romanzo,
          mai nei tweet.
    \item \textit{had been} (past perfect) è molto frequente nel romanzo
          (18,85\%), assente nei tweet.
\end{itemize}

\subsubsection{Limite del MLE: il problema dello zero}

Il principale limite del MLE è che assegna probabilità \textbf{zero}
a qualsiasi sequenza non osservata nel training. Questo causa la
\textbf{perplexity infinita} quando il testo di test contiene parole
o bigrammi mai visti:

\begin{equation}
P_{\text{MLE}}(\text{``xyzword''} \mid \text{``i''}) = \frac{0}{C(\text{``i''})} = 0
\end{equation}

\begin{lstlisting}[language=Python, caption=Dimostrazione del problema dello zero]
print(modello_tweet.score('xyzword', ['i']))
# -> 0.000000  (zero esatto)

frase_rara = word_tokenize("i xyzword lol")
print(modello_tweet.perplexity(frase_rara))
# -> inf  (infinita!)
\end{lstlisting}

\subsection{Smoothing di Laplace}

Il \textbf{Laplace smoothing} (o add-one smoothing) risolve il problema
dello zero aggiungendo 1 a tutti i conteggi:

\begin{equation}
P_{\text{Laplace}}(w_i \mid w_{i-1}) =
\frac{C(w_{i-1}, w_i) + 1}{C(w_{i-1}) + |V|}
\end{equation}

dove $|V|$ è la dimensione del vocabolario.

\begin{lstlisting}[language=Python, caption=Addestramento modello Laplace]
from nltk.lm import Laplace

modello_lp = Laplace(n)
modello_lp.fit(train_data, vocab)

# Ora le parole mai viste hanno probabilita' piccola ma > 0
print(modello_lp.score('xyzword', ['i']))
# -> 0.000065  (piccolo ma non zero)

print(modello_lp.perplexity(frase_rara))
# -> 1677.95  (alta ma non infinita)
\end{lstlisting}

\begin{table}[H]
\centering
\begin{tabular}{lrr}
\toprule
\textbf{Metrica} & \textbf{MLE} & \textbf{Laplace} \\
\midrule
P('xyzword' $|$ 'i')          & 0,000000 & 0,000065 \\
Perplexity su parola rara     & $\infty$ & 1677,95 \\
Gestisce parole mai viste     & No       & Sì \\
\bottomrule
\end{tabular}
\caption{Confronto MLE vs Laplace}
\end{table}

\subsection{Scelta del valore di N}

\begin{table}[H]
\centering
\begin{tabular}{lrr}
\toprule
\textbf{N} & \textbf{PP frase Twitter} & \textbf{PP frase Letterario} \\
\midrule
1 & 5170,14 & 12025,00 \\
2 & 5170,91 & 12027,00 \\
3 & 5170,92 & 12027,00 \\
\bottomrule
\end{tabular}
\caption{Confronto perplexity per n=1,2,3 (Laplace, modello Twitter)}
\end{table}

I valori quasi identici per $N=1,2,3$ indicano un problema di
\textbf{data sparsity}: con soli 94.706 token Twitter, la maggior
parte dei bigrammi e trigrammi appare raramente, e Laplace li tratta
quasi tutti come sequenze rare. Con un corpus più grande le differenze
sarebbero più marcate, e i trigrammi catturerebbero strutture
sintattiche più complesse.

% ============================================================
\section{Perplexity e Classificazione}
% ============================================================

\subsection{Cos'è la Perplexity}

La \textbf{perplexity} misura quanto un modello è ``sorpreso'' da un testo:

\begin{equation}
\text{PP}(W) = P(w_1, w_2, \ldots, w_N)^{-\frac{1}{N}} =
\sqrt[N]{\frac{1}{P(w_1, w_2, \ldots, w_N)}}
\end{equation}

Un valore \textbf{basso} indica che il modello conosce bene quel tipo
di testo; un valore \textbf{alto} indica sorpresa.

\subsection{Perplexity cross-domain}

\begin{table}[H]
\centering
\begin{tabular}{llrrl}
\toprule
\textbf{Frase} & \textbf{Tipo atteso} &
\textbf{PP Twitter} & \textbf{PP Lett.} & \textbf{Verdetto} \\
\midrule
``i love you so much''       & TW & 5170,9 & 4488,5 & TW \checkmark \\
``so happy right now''       & TW & 12027,0 & 9308,0 & LT $\times$ \\
``i miss you feel so sad''   & TW & 5746,3 & 4917,0 & TW \checkmark \\
``she had been very well''   & LT & 12027,0 & 9308,0 & TW $\times$ \\
``it was a very good thing'' & LT & 6756,2 & 5258,0 & TW $\times$ \\
\bottomrule
\end{tabular}
\caption{Perplexity cross-domain (Laplace normalizzata)}
\end{table}

\textbf{Limite osservato}: la perplexity con Laplace su vocabolari
di dimensioni diverse non è direttamente comparabile. Il modello con
vocabolario più grande (Twitter, 12.027 parole) distribuisce la
probabilità su più termini e ottiene sempre perplexity più alta,
indipendentemente dal contenuto del testo.

\subsection{Classificatore con Log-Likelihood}

Per superare il limite della perplexity, abbiamo implementato un
classificatore basato sulla \textbf{log-likelihood}:

\begin{equation}
\log P(w_1, \ldots, w_N) = \sum_{i=2}^{N} \log P(w_i \mid w_{i-1})
\end{equation}

Un valore \textbf{meno negativo} (più alto) indica che il modello
spiega meglio il testo.

\begin{lstlisting}[language=Python, caption=Funzione log-likelihood]
import math

def log_likelihood(modello, tokens):
    score = 0.0
    for i in range(1, len(tokens)):
        contesto = [tokens[i-1]]
        prob = modello.score(tokens[i], contesto)
        if prob > 0:
            score += math.log(prob)
        else:
            score += math.log(1e-10)  # parole mai viste
    return score

def classifica_loglik(testo, mod_tw, mod_lt):
    tokens = word_tokenize(testo.lower())
    ll_tw   = log_likelihood(mod_tw, tokens)
    ll_lett = log_likelihood(mod_lt, tokens)
    return "TWITTER" if ll_tw > ll_lett else "LETTERARIO"
\end{lstlisting}

\subsubsection{Risultati del classificatore}

\begin{table}[H]
\centering
\begin{tabular}{p{7cm}rrr}
\toprule
\textbf{Testo} &
\textbf{LL Twitter} & \textbf{LL Lett.} & \textbf{Verdetto} \\
\midrule
``i love you so much thank you''
    & $-32,02$ & $-42,15$ & TWITTER \checkmark \\
``she had been very well disposed''
    & $-124,74$ & $-21,91$ & LETTERARIO \checkmark \\
``i miss you so much cant stop thinking''
    & $-41,54$ & $-122,88$ & TWITTER \checkmark \\
``it was a very good thing to have done''
    & $-29,21$ & $-24,75$ & LETTERARIO \checkmark \\
``happy birthday love you so much''
    & $-11,27$ & $-58,54$ & TWITTER \checkmark \\
``she was not a woman of many words''
    & $-70,82$ & $-47,21$ & LETTERARIO \checkmark \\
``i cant stop smiling so happy''
    & $-57,95$ & $-97,56$ & TWITTER \checkmark \\
``mr knightley had always been very fond''
    & $-108,42$ & $-24,91$ & LETTERARIO \checkmark \\
\bottomrule
\end{tabular}
\caption{Risultati del classificatore log-likelihood (8/8 corretti)}
\end{table}

Il classificatore log-likelihood classifica correttamente tutti gli 8 testi di test,
a differenza del classificatore basato su perplexity (3/8).

% ============================================================
\section{Generazione del Testo}
% ============================================================

Il modello MLE può essere usato per \textbf{generare} testo nello
stile del corpus su cui è stato addestrato. La generazione avviene
campionando la parola successiva in base alle probabilità condizionali
apprese.

\begin{lstlisting}[language=Python, caption=Generazione testo]
def genera_testo(modello, n_parole=30):
    testo_raw = modello.generate(n_parole, random_seed=42)
    # Filtra token speciali e lettere singole (artefatti)
    testo_filtrato = [
        t for t in testo_raw
        if t not in ['<s>', '</s>', '<UNK>']
        and not (len(t) == 1 and t not in ['i', 'a'])
    ]
    return ' '.join(testo_filtrato)
\end{lstlisting}

\subsubsection{Testo generato}

\begin{table}[H]
\centering
\begin{tabular}{p{2cm}p{10cm}}
\toprule
\textbf{Stile} & \textbf{Testo generato} \\
\midrule
Twitter &
``papa coach it from the rain why are not a fantastic spring
fiesta make sure the donation manager we all the rest his cute
when im being a small polka dots'' \\
\midrule
Letterario &
``over a good exchange consent to the cardroom looking about
her less allied with persevering attention could not that a
substance of her but yet heard any body to me'' \\
\bottomrule
\end{tabular}
\caption{Esempi di testo generato dai due modelli}
\end{table}

La differenza stilistica è evidente: il testo Twitter è disorganizzato
e colloquiale, mentre quello letterario usa costruzioni più complesse
tipiche del romanzo ottocentesco.

\subsubsection{Confronto con diversi seed}

\begin{table}[H]
\centering
\begin{tabular}{cp{5cm}p{5cm}}
\toprule
\textbf{Seed} & \textbf{Twitter} & \textbf{Letterario} \\
\midrule
99  & ``im getting bigger at the epic i love you in my life'' &
      ``i can believe i shall decidedly in the neighbourhood'' \\
7   & ``have a new cards in march free insurance'' &
      ``hartfield and return among their getting at last'' \\
123 & ``and by my bestfriend she brought us i want corn chips'' &
      ``and away from being to announce cheerful exulting'' \\
\bottomrule
\end{tabular}
\caption{Generazione con diversi seed casuali}
\end{table}

% ============================================================
\section{Analisi Linguistica}
% ============================================================

\subsection{Type-Token Ratio (TTR)}

Il \textbf{Type-Token Ratio} misura la ricchezza lessicale di un corpus:

\begin{equation}
\text{TTR} = \frac{\text{numero di parole uniche (types)}}
                  {\text{numero totale di parole (tokens)}}
\end{equation}

\begin{table}[H]
\centering
\begin{tabular}{lrrl}
\toprule
\textbf{Corpus} & \textbf{Types} & \textbf{Tokens} & \textbf{TTR} \\
\midrule
Twitter    & 12.024 & 94.706  & 0,1270 \\
Letterario & 9.305  & 158.267 & 0,0588 \\
\bottomrule
\end{tabular}
\caption{Type-Token Ratio dei due corpora}
\end{table}

Il TTR di Twitter è più del \textbf{doppio} rispetto al romanzo.
Questo risultato può sembrare controintuitivo, ma si spiega con
la natura eterogenea dei tweet: ogni tweet introduce nomi propri,
hashtag, abbreviazioni e slang diversi (\texttt{followfriday},
\texttt{unfollowers}, \texttt{bestfriend}), molti dei quali appaiono
raramente nel corpus. Il romanzo invece ripete sistematicamente gli
stessi personaggi e costruzioni grammaticali.

\subsection{Lunghezza media delle frasi}

\begin{table}[H]
\centering
\begin{tabular}{lr}
\toprule
\textbf{Corpus} & \textbf{Lunghezza media (parole)} \\
\midrule
Tweet             & 16,0 \\
Frase letteraria  & 25,6 \\
\bottomrule
\end{tabular}
\caption{Lunghezza media delle frasi}
\end{table}

Le frasi letterarie sono in media \textbf{9,6 parole più lunghe} dei
tweet. Questo riflette la struttura sintattica complessa del romanzo
ottocentesco, con subordinate multiple, incisi e descrizioni elaborate,
contrapposta alla brevità imposta dal limite di 140 caratteri di Twitter.

% ============================================================
\section{Limiti e Vantaggi}
% ============================================================

\subsection{Vantaggi dell'approccio N-gram}

\begin{itemize}
    \item \textbf{Semplicità}: facile da implementare e interpretare.
    \item \textbf{Efficienza}: addestramento e inferenza molto veloci.
    \item \textbf{Trasparenza}: le probabilità condizionali sono
          direttamente leggibili e spiegabili.
    \item \textbf{Nessun training pesante}: non richiede GPU o
          risorse computazionali elevate.
\end{itemize}

\subsection{Limiti dell'approccio N-gram}

\begin{itemize}
    \item \textbf{Data sparsity}: con dataset piccoli (come i 10.000
          tweet usati), la maggior parte dei bigrammi e trigrammi
          appare raramente o mai. Questo rende i valori di perplexity
          per $N=1,2,3$ quasi identici.

    \item \textbf{Finestra di contesto limitata}: i bigrammi
          considerano solo la parola precedente. Non catturano
          dipendenze a lungo raggio tipiche del linguaggio naturale
          (``The man who I met yesterday \textbf{was} kind'').

    \item \textbf{Nessuna comprensione semantica}: il modello non
          capisce il significato delle parole. \textit{happy} e
          \textit{joyful} sono trattate come parole completamente
          diverse.

    \item \textbf{Problema della perplexity con Laplace}:
          la perplexity non è direttamente comparabile tra modelli
          con vocabolari di dimensioni diverse, rendendo necessario
          l'uso della log-likelihood per il classificatore.

    \item \textbf{Dataset sbilanciato}: il corpus Twitter (94.706
          token) è significativamente più piccolo del corpus
          letterario (158.267 token), influenzando il confronto.
\end{itemize}

\subsection{Confronto con approcci alternativi}

\begin{table}[H]
\centering
\begin{tabular}{p{3cm}p{3.5cm}p{3.5cm}p{2.5cm}}
\toprule
\textbf{Approccio} & \textbf{Vantaggi} &
\textbf{Svantaggi} & \textbf{Complessità} \\
\midrule
N-gram (questo lab) &
Semplice, veloce, interpretabile &
Nessun contesto semantico, sparsità &
Bassa \\
\midrule
Word2Vec &
Cattura semantica, vettori densi &
Non genera testo, no probabilità &
Media \\
\midrule
LSTM / RNN &
Contesto lungo, genera testo realistico &
Richiede molto training &
Alta \\
\midrule
Transformer (GPT) &
Stato dell'arte, testo coerente &
Computazionalmente costoso &
Molto alta \\
\bottomrule
\end{tabular}
\caption{Confronto tra approcci per la modellazione del linguaggio}
\end{table}

% ============================================================
\section{Risultati e Conclusioni}
% ============================================================

\subsection{Riepilogo dei risultati}

\begin{table}[H]
\centering
\begin{tabular}{lll}
\toprule
\textbf{Esperimento} & \textbf{Risultato} & \textbf{Interpretazione} \\
\midrule
Bigrammi Twitter     & (i, love), (cant, wait) &
Linguaggio emozionale, informale \\
Bigrammi Letterario  & (mr, knightley), (mrs, weston) &
Nomi propri formali, titoli \\
TTR Twitter          & 0,1270 &
Vocabolario più vario \\
TTR Letterario       & 0,0588 &
Vocabolario più ripetitivo \\
Lunghezza tweet      & 16,0 parole &
Frasi brevi \\
Lunghezza frase lett. & 25,6 parole &
Frasi lunghe e complesse \\
Classificatore (PP)  & 3/8 = 38\% &
Inaffidabile \\
Classificatore (LL)  & 8/8 = 100\% &
Accurato su frasi note \\
\bottomrule
\end{tabular}
\caption{Riepilogo dei principali risultati}
\end{table}

\subsection{Conclusioni}

Questo laboratorio ha dimostrato che i modelli N-gram sono strumenti
efficaci per caratterizzare e distinguere stili linguistici diversi.

I principali risultati ottenuti sono:

\begin{enumerate}
    \item I bigrammi e trigrammi più frequenti riflettono chiaramente
          le differenze culturali e stilistiche tra Twitter e il
          linguaggio letterario ottocentesco.

    \item Il TTR più alto di Twitter conferma che i social media
          producono un linguaggio lessicalmente più vario,
          sebbene meno strutturato sintatticamente.

    \item La perplexity con Laplace non è un indicatore affidabile
          per la classificazione cross-dominio quando i vocabolari
          hanno dimensioni molto diverse. La log-likelihood è
          una metrica più robusta per questo scopo.

    \item Il classificatore basato su log-likelihood raggiunge
          un'accuratezza del 100\% sulle frasi di test, confermando
          che i due domini sono statisticamente ben separabili.

    \item La generazione di testo evidenzia le differenze stilistiche:
          il modello Twitter produce testo colloquiale e frammentato,
          mentre il modello letterario produce costruzioni più
          elaborate tipiche del romanzo.
\end{enumerate}

Sviluppi futuri potrebbero includere l'uso di dataset più grandi
per ridurre la sparsità, l'adozione di smoothing più sofisticati
come Kneser-Ney, e il confronto con modelli neurali come LSTM o
Transformer per valutare il guadagno in qualità rispetto alla
complessità computazionale.

% ============================================================
\section*{Contributi da strumenti generativi}
% ============================================================

\addcontentsline{toc}{section}{Contributi da strumenti generativi}

Come richiesto dalle istruzioni del laboratorio, si dichiara che
nella stesura di questa relazione e nella scrittura del codice
sono stati utilizzati strumenti di AI generativa (Claude, Anthropic)
per:

\begin{itemize}
    \item Supporto nella strutturazione del codice Python.
    \item Debugging di errori nella generazione del testo e nel
          classificatore.
    \item Revisione e formattazione della relazione in \LaTeX.
\end{itemize}

Tutte le scelte progettuali, l'analisi dei risultati e le conclusioni
sono state elaborate autonomamente dallo studente.

% ============================================================
\end{document}
% ============================================================